\documentclass{article}
\usepackage[utf8]{inputenc}
\usepackage[T1]{fontenc}
\usepackage{geometry}
\usepackage{amsmath}
\usepackage{amssymb}
\usepackage{graphicx}
\usepackage{xcolor}
\usepackage{hyperref}
\usepackage{titlesec}
\usepackage{array}

% Page setup
\geometry{a4paper, margin=1in}
\setlength{\parindent}{0pt}
\setlength{\parskip}{1em}

% Section styling
\titleformat{\section}
  {\normalfont\Large\bfseries\color{blue!70!black}}
  {\thesection}{1em}{}

\titleformat{\subsection}
  {\normalfont\large\bfseries\color{blue!50!black}}
  {\thesubsection}{1em}{}

% Table styling
\newcolumntype{P}[1]{>{\raggedright\arraybackslash}p{#1}}

\begin{document}

\begin{titlepage}
    \centering
    \vspace*{2cm}
    {\Huge\bfseries The Hidden Heritage of Manipur\par}
    \vspace{1cm}
    {\Large Culture, Manuscripts, and Language\par}
    \vspace{2cm}
    {\large A glimpse into the rich traditions of the “Jewel of India”\par}
    \vspace{3cm}
    {\today\par}
\end{titlepage}

\tableofcontents
\newpage

\section{Vibrant Culture of Manipur}
Manipur, a northeastern state of India, is often called the “Jewel of India” for its breathtaking landscape and its rich, distinct cultural heritage. Despite being a part of India, its culture remains largely unfamiliar to the majority of the Indian population due to geographical isolation and historical complexities.

\subsection{Classical Dance: Ras Leela}
The \textbf{Manipuri Ras Leela} is one of the eight classical dance forms of India. It is deeply devotional, depicting the divine love between Radha and Krishna. Unlike other classical dances, it is characterized by graceful, fluid movements, delicate footwork, and the use of the \textit{pung} (a classical drum). The dance is not merely performance but a form of worship, often staged within the sacred precincts of temples like the Shri Shri Govindajee Temple in Imphal.

\subsection{Traditional Festivals}
Manipur’s calendar is dotted with festivals that reflect its agrarian roots and spiritual depth.
\begin{itemize}
    \item \textbf{Yaoshang}: Often compared to Holi, Yaoshang is a five-day spring festival featuring traditional sports, theatrical performances, and the distinctive \textit{Thabal Chongba} (moonlight dancing).
    \item \textbf{Lai Haraoba}: This is the oldest festival, meaning “merrymaking of the gods.” It celebrates the creation myths and the pre-Hindu Sanamahi religion through elaborate rituals, songs, and dances that preserve the state’s ancient folklore.
    \item \textbf{Ningol Chakkouba}: A unique social festival where married women are invited back to their parental homes for a grand feast, reinforcing familial bonds.
\end{itemize}

\subsection{Traditional Attire and Martial Arts}
The \textbf{Phanek} (a sarong-like wrap) and \textbf{Innaphi} (a shawl) form the traditional attire, known for their distinctive geometric patterns. Manipur is also the birthplace of \textbf{Thang-Ta}, a sophisticated martial art form that incorporates the spear (\textit{Thang}) and sword (\textit{Ta}) with intricate footwork and ritualistic elements. This art form is not just combat training but a philosophical discipline.

\begin{center}
    \begin{tabular}{|P{4cm}|P{8cm}|}
        \hline
        \textbf{Cultural Element} & \textbf{Significance} \\
        \hline
        Ras Leela & Classical dance, spiritual storytelling \\
        \hline
        Lai Haraoba & Preservation of pre-Vedic rituals and cosmology \\
        \hline
        Thang-Ta & Martial art embodying physical and spiritual discipline \\
        \hline
    \end{tabular}
    \captionof{table}{Key cultural elements of Manipur}
\end{center}

\section{Manuscripts of Manipur}
Manipur possesses a vast and largely hidden corpus of manuscripts that preserve its history, philosophy, and science. These texts, written on materials like palm leaves (\textit{Matam}) and handmade paper, constitute a unique intellectual heritage.

\subsection{Puyas: The Ancient Texts}
The \textbf{Puyas} are a collection of ancient Manipuri texts that document the history, genealogy of kings, and the philosophical tenets of the Sanamahi religion. Some of these texts, such as the \textit{Wakoklon Thilel Salai Amailon Pukok Puya}, are considered to predate the Hindu epics and contain detailed accounts of cosmology, ethics, and statecraft. For centuries, these manuscripts were preserved in royal archives and family libraries, but many were lost or destroyed during colonial rule and the ensuing historical upheavals.

\subsection{The Kangla and Archival Collections}
The \textbf{Kangla}, the ancient capital of Manipur, was the center of manuscript preservation. Today, institutions like the \textbf{Manipur State Archives} and the \textbf{Indira Gandhi National Centre for the Arts (IGNCA)} regional center in Imphal have undertaken digitization projects to preserve these fragile texts. The manuscripts cover diverse subjects:
\begin{itemize}
    \item \textbf{Astronomy and Astrology}: Detailed charts and calculations.
    \item \textbf{Medicine}: Traditional healing practices (\textit{Mai-sa}) using local herbs.
    \item \textbf{Literature}: Epic poetry and chronicles (\textit{Cheitharol Kumbaba}), the royal chronicle of Manipur.
\end{itemize}

\subsection{Preservation Challenges}
Despite their immense value, these manuscripts face threats from humidity, insects, and lack of adequate conservation infrastructure. Many remain in private hands, inaccessible to scholars. The effort to catalog and digitize them is ongoing but remains underfunded and under-recognized at the national level.

\begin{center}
    \fbox{\parbox{0.9\textwidth}{\centering
    \textbf{Note:} The \textit{Cheitharol Kumbaba} is one of the most important historical manuscripts, chronicling the reigns of Manipuri kings from 33 CE to 1955 CE, making it one of the longest continuous royal chronicles in the world.
    }}
\end{center}

\section{Local Languages of Manipur}
Linguistic diversity in Manipur is complex and often overshadowed by the dominance of Meitei (Manipuri). The state is a linguistic microcosm with over 30 languages belonging to the Tibeto-Burman family.

\subsection{Meiteilon (Manipuri Language)}
Meiteilon, or Manipuri, is the official language of the state and was recognized as a “Classical Language” of India in 2024 due to its ancient lineage and rich literary tradition. It has its own distinct script, the \textbf{Meitei Mayek}, which was in use for centuries before the adoption of the Bengali script during the 18th century. There is a contemporary movement to revive the Meitei Mayek script in education and public life.

\subsection{Tribal Languages}
The hill districts of Manipur are home to numerous indigenous tribal communities, each with its own language. Key languages include:
\begin{itemize}
    \item \textbf{Tangkhul}: Spoken in Ukhrul district, with numerous dialects.
    \item \textbf{Thadou (Kuki)}: A Kuki-Chin language widely spoken in the southern hills.
    \item \textbf{Kabui (Rongmei)}: Part of the Zeliangrong community.
    \item \textbf{Paite, Hmar, Vaiphei, and Zou}: Other significant languages belonging to the Kuki-Chin subgroup.
\end{itemize}
These languages are largely oral, though missionary efforts and community initiatives have developed written scripts, primarily using the Roman alphabet.

\subsection{Linguistic Marginalization}
Despite constitutional protections, many of these local languages face existential threats. The dominance of Meitei and English in education, governance, and media, combined with historical ethnic conflicts, has led to language shift and endangerment. Organizations like the \textbf{Tribal Literature Forum} work to document and promote these languages through literature and education.

\subsection{Comparative Language Table}
\begin{center}
    \begin{tabular}{|l|l|l|}
        \hline
        \textbf{Language} & \textbf{Language Family} & \textbf{Status} \\
        \hline
        Meitei (Manipuri) & Tibeto-Burman & Classical Language, Official \\
        \hline
        Tangkhul & Tibeto-Burman & Scheduled Language (in hill districts) \\
        \hline
        Thadou & Tibeto-Burman & Widely spoken, recognized \\
        \hline
        Rongmei & Tibeto-Burman & Endangered \\
        \hline
        Anal & Tibeto-Burman & Vulnerable \\
        \hline
    \end{tabular}
    \captionof{table}{Major languages of Manipur}
\end{center}

\subsection*{References}
\begin{thebibliography}{9}
    \bibitem{singh} Singh, R.K. Jhalajit. \textit{A Short History of Manipur}. 4th ed., Manipur: RKJ, 2014.
    \bibitem{parratt} Parratt, Saroj Nalini Arambam. \textit{The Court Chronicle of the Kings of Manipur: The Cheitharon Kumpapa}. Routledge, 2005.
    \bibitem{devi} Devi, L. Kunjeswori. \textit{Archaeology in Manipur}. Akansha Publishing House, 2006.
    \bibitem{ignca} Indira Gandhi National Centre for the Arts. \textit{Manuscripts of Manipur: A Digital Collection}. IGNCA Regional Centre, Imphal, 2019.
    \bibitem{chelliah} Chelliah, Shobhana L. \textit{A Grammar of Meithei}. Mouton de Gruyter, 1997.
    \bibitem{shakespeare} Shakespeare, John. \textit{The Kuki-Lushai Clans}. Macmillan, 1912 (reprint: Tribal Research Institute, Manipur).
    \bibitem{government} Government of Manipur. \textit{Gazetteer of Manipur}. Government Press, Imphal, 2020.
\end{thebibliography}

\end{document}